\documentclass[11pt]{article}
\usepackage[margin=1in]{geometry}
\usepackage{amsmath,amssymb,amsthm,booktabs,longtable}
\usepackage[colorlinks,linkcolor=blue,citecolor=blue]{hyperref}

\theoremstyle{plain}
\newtheorem{theorem}{Theorem}
\newtheorem{proposition}[theorem]{Proposition}
\newtheorem{lemma}[theorem]{Lemma}
\newtheorem{corollary}[theorem]{Corollary}
\theoremstyle{definition}
\newtheorem{remark}[theorem]{Remark}
\newtheorem{measurement}[theorem]{Measurement}

\newcommand{\edot}{\dot{e}}
\newcommand{\Par}{\mathrm{Par}}

\title{Q436 and Q435: the branch assignment, and what $t=0$ actually counts\\
{\large\normalfont Warnaar's $t$ grades the winding number, not a Hall--Littlewood parameter}}
\author{Clio}
\date{13 October 2026 (PROVE, cycle 1)}

\begin{document}
\maketitle

\begin{center}
\fbox{\begin{minipage}{0.92\textwidth}
\small
\textbf{Recipient:} Rick (\texttt{grandparick20@gmail.com}), cc Robin. \quad
\textbf{Author:} Clio. \quad \textbf{Date:} 13 October 2026.\\
\textbf{Source commit:} \texttt{clio-vega/proofs@224e14a} --- the commit that first published
this \texttt{.tex}, and the content this PDF was built from. Per PROTOCOL \S2.3; the hash names
the source, so the commit carrying \emph{this} line is necessarily its child.\\
\textbf{Registry:} \texttt{proofs/registry/cylindric-statistic.json}. Six nodes cite this file:
\texttt{t0-selects-zero-winding} (\emph{proved}),
\texttt{c-at-zero-is-unwrapped-count} (\emph{computed}),
\texttt{both-branches-multiplicity-determined} (\emph{proved}),
\texttt{palazzo-does-not-determine-branch} (\emph{proved}),
\texttt{korff-f-vs-N-conflation} (\emph{proved}),
\texttt{morse-schilling-index-set-mismatch} (\emph{computed}).
\end{minipage}}
\end{center}

\begin{abstract}
The session was briefed as a read plus arithmetic: assign the Macdonald branch of the
target $c_{(1^4)}(t)=3-t^2$ (Q436), then test the hypothesis that $c_\alpha(0)$ is an
undeformed coset count (Q435). Both questions are answered, and both answers correct the
brief. (i)~Korff--Palazzo \texttt{1804.05647} does \emph{not} determine the branch: it names
the $q$-Whittaker specialisation once, without a display, and its own letter $t$ is a third
parameter. The primary source Korff \texttt{1110.6356}~\S2.6 gives both specialisations
explicitly, and his $f^\nu_{\lambda\mu}(t)$ is Hall--Littlewood. (ii)~The branch assignment
is nevertheless \emph{irrelevant to every weight kill}, because both specialisations have
multiplicity-determined single-box skew weights and are therefore both killed by Theorem~D of
2026-10-12 (Theorem~\ref{thm:bothbranches} below, $280/280$). (iii)~The real finding is
Theorem~\ref{thm:t0}: in Warnaar's $\mathrm{Eq}_{\mathrm{generalise}}$ every $t$-exponent is
non-negative and vanishes only at zero winding, so $t=0$ selects the \emph{unwrapped}
determinant. Hence $c_\alpha(0)$ is the count of zero-winding tableaux
(Theorem~\ref{thm:count}, $1656/1656$), the undeformed \emph{cylindric} count is $c_\alpha(1)$,
and the briefed hypothesis is refuted on its pre-registered arm~A2 with the discrepancy
located at the normalisation end. Warnaar's $t$ sits in Korff's winding slot $z$, not in his
Hall--Littlewood slot $t$ nor in his $q$-Whittaker slot $q$ --- which is a structural
explanation of why four successive per-box weight architectures died.
\end{abstract}

\section*{\S0 \ Holdings, measured by this session}

\textbf{Root list searched, printed because a holdings verdict's domain \emph{is} its root
list:} \texttt{/home/clio/data}, \texttt{/home/clio/projects}, \texttt{/home/clio/git},
\texttt{/home/clio/mail}. Every file below was \texttt{stat}-ed by this session at
13 October, start of slot; sizes are bytes as measured, not copied from the brief.

\begin{center}\small
\begin{tabular}{@{}llr@{}}
\toprule
paper & path measured present & bytes \\
\midrule
Korff--Palazzo \texttt{1804.05647} & \texttt{projects/sources/20260924c1-novelty/korff-palazzo-1804.05647.tex} & $235{,}421$ \\
Morse--Schilling \texttt{1207.0786} & \texttt{projects/library/.../wk/ms-1207.0786.pdf} & $109{,}844$ \\
Korff \texttt{1110.6356} (pdf) & \texttt{projects/library/.../wk/korff-1110.6356.pdf} & $1{,}126{,}690$ \\
Korff \texttt{1110.6356} (text layer) & \texttt{projects/library/.../wk/korff-1110.6356.txt} & $202{,}046$ \\
HKKO \texttt{2301.13117} & \texttt{/home/clio/data/fetched/arxiv-2301.13117.pdf} & $662{,}217$ \\
Wheeler--Zinn-Justin \texttt{1603.01815} & \texttt{git/puzzles/seed-papers/wheeler-zinn-justin-2016-\dots pdf} & $665{,}583$ \\
\bottomrule
\end{tabular}
\end{center}

All six present. Line numbers quoted below refer to the \texttt{1110.6356} text layer
(\texttt{korff-1110.6356.txt}, $7{,}650$ lines, $38{,}506$ words) and to the
\texttt{1804.05647} \LaTeX{} source ($2{,}853$ lines), both read first-hand this session.
Morse--Schilling was re-extracted this session with \texttt{pdftotext}
($4{,}740$ words, $700$ lines).

\section{\S1 \ Q436: the branch assignment}

\subsection{What Korff--Palazzo \texttt{1804.05647} does and does not say}

\begin{measurement}\label{m:palazzo}
\texttt{1804.05647} is Korff--Palazzo, \emph{Cylindric symmetric functions and positivity}
(title, l.\,132). The string \texttt{Whittaker} occurs \textbf{exactly once} in the whole
$25{,}851$-word source (\texttt{grep -ci whittaker} $=1$), at l.\,182, in the
introduction, describing a \emph{different} paper:
\begin{quote}\small
``\dots a $q$-deformation of the $\widehat{sl}_n$-Verlinde algebra was constructed in
[korff2013cylindric] which \dots whose structure constants or fusion coefficients
$N^\nu_{\lambda\mu}(q)\in\mathbb{Z}[q]$ are related to cylindric versions of
Hall--Littlewood and $q$-Whittaker polynomials. Both types of polynomials occur (in the
non-cylindric case) as specialisation of Macdonald polynomials.'' (l.\,182)
\end{quote}
No display, no substitution into $P_\lambda(x;q,t)$. The cited
\texttt{korff2013cylindric} is, from the bibliography (ll.\,2705--2709), Korff,
\emph{Cylindric versions of specialised Macdonald functions and a deformed Verlinde
algebra}, CMP \textbf{318} (2013) 173--246 $=$ arXiv \texttt{1110.6356}.
\end{measurement}

\begin{measurement}[answering \S1(iv)]\label{m:thirdparam}
\emph{Palazzo's $t$ is a third parameter.} In \texttt{1804.05647} the letter $t$ is a
cyclotomic Hecke/twist variable, not a Macdonald parameter: l.\,429 defines the deformed
$\mathcal{S}$-matrix $\mathcal{S}(t)=\frac{1}{\sqrt n}\sum_{i,j}t^{-j}\zeta^{ij}e_{ij}$ with
$\zeta=\exp(2\pi\iota/n)$ (l.\,334); the lemma after it sets $z=t^n$; and l.\,684 is the
cyclotomic polynomial $f(x_1)=(x_1-t)(x_1-t\zeta)\cdots(x_1-t\zeta^{n-1})=x_1^n-t^n$. At
l.\,953 this $t$ multiplies the \emph{undeformed} fusion coefficient as a pure grading
factor, $\langle v^\nu, M^*_\lambda v_\mu\rangle = t^{|\nu|-|\lambda|-|\mu|}N^\nu_{\lambda\mu}$,
with $N^\nu_{\lambda\mu}=0$ unless $|\lambda|+|\mu|-|\nu|\equiv 0 \bmod n$ (l.\,955). The
body of the paper works at $q\to 1$ (l.\,184: ``we investigate the combinatorial structure
of the $q$-deformed Verlinde algebra from [korff2013cylindric] \textbf{in the limit
$q\to 1$}'').
\end{measurement}

\begin{corollary}[\S1(i), answered in the negative]
Korff--Palazzo \texttt{1804.05647} does not determine the branch of any $t$-dependent target.
It mentions the two specialisations once without displaying either, works only at $q\to 1$,
and uses the letter $t$ for an unrelated cyclotomic parameter. \emph{This is a result, not a
failure of the read} --- the brief named the wrong paper, and the right one is on disk.
\end{corollary}

\subsection{The two specialisations, from the primary source}

\begin{measurement}[\S1(i) from \texttt{1110.6356}~\S2.6, read first-hand]\label{m:twobranches}
Both specialisations are displayed in Korff \texttt{1110.6356}~\S2.6.
\begin{itemize}
\item \textbf{Hall--Littlewood branch: $q=0$, surviving parameter $t$.} ``In what follows we
will omit the dependence on the second parameter in the notation and simply write
$Q_\lambda(x;t)=Q_\lambda(x;0,t)$ and $P_\lambda(x;t)=P_\lambda(x;0,t)$'' (l.\,1313). The
skew weights are, eq.~(2.21),
\[
  \psi_{\lambda/\mu}(t)=\prod_{i\in J_{\lambda/\mu}}\bigl(1-t^{\mu'_i-\mu'_{i+1}}\bigr),
  \qquad J_{\lambda/\mu}=\{\,i:\theta'_i=0,\ \theta'_{i+1}=1\,\},\ \ \theta'=\lambda'-\mu',
\]
zero unless $\lambda/\mu$ is a horizontal strip, with $m_i(\mu)=\mu'_i-\mu'_{i+1}$ the
multiplicity of the part $i$ (l.\,1296).
\item \textbf{$q$-Whittaker branch: $t=0$, surviving parameter $q$.} ``In light of the
definition (2.19) it is natural to consider also the complementary limit of Macdonald
polynomials setting now $t=0$'' (l.\,1349). Eq.~(2.24):
\[
  \psi'_{\lambda/\mu}(q):=\psi_{\lambda/\mu}(q,0)
  =\prod_{i>0}\begin{bmatrix}\lambda_i-\lambda_{i+1}\\ \lambda_i-\mu_i\end{bmatrix}_q,
\]
again zero off horizontal strips, with the companion $\varphi'_{\lambda/\mu}(q)$ at
eq.~(2.25). Korff adds (l.\,1427): ``In [26] they have been named $q$-deformed or simply
$q$-Whittaker functions.''
\end{itemize}
\end{measurement}

\begin{measurement}[\S1(ii)]\label{m:fisHL}
\emph{Korff's $f^\nu_{\lambda\mu}(t)$ is on the Hall--Littlewood branch.} It is introduced
in \S2.6.2, titled ``Hall algebra, coproduct and skew Hall--Littlewood functions'', as the
coefficients in the product expansion (1.3) of Hall--Littlewood $P$-functions, and
Example~7.2 calls them ``the expansion coefficients in the corresponding product of
Hall--Littlewood polynomials'' (l.\,7230). The $q$-Whittaker objects carry primes
($\psi'$, $\varphi'$, $P'$, $Q'$) and the letter $q$ throughout. There is no ambiguity.
\end{measurement}

\begin{measurement}[the brief's control, corrected]\label{m:fvsN}
\textbf{$f$ and $N$ are different objects, and the brief's ``undeformed fusion/coset count''
is a Littlewood--Richardson coefficient.} In Example~7.2 ($n=4$, $k=3$, $z=1$,
$\lambda=(3,2,1)$, $\mu=(4,3,1)$) Korff writes, l.\,7177:
\begin{quote}\small
``We exploit the fact that $f^\nu_{\lambda\mu}(t)\equiv 0$ unless
$f^\nu_{\lambda\mu}(0)=c^\nu_{\lambda\mu}\neq 0$ \textbf{where $c^\nu_{\lambda\mu}$ is the
Littlewood--Richardson coefficient}.''
\end{quote}
So the briefed ``$c^{(6,5,3)}_{\lambda\mu}=2$'' is an LR coefficient. The
\emph{deformed fusion} coefficients are the separate family
$N^\lambda_{\mu\nu}(t)$ obtained from the $f$'s by signed straightening
(ll.\,7256--7262), and for them Korff states, l.\,7266: ``Setting $t=0$ one can verify,
using the Verlinde formula or the Kac--Walton algorithm, that one obtains the correct fusion
coefficients of the $\widehat{sl(n)}_k$-Verlinde algebra.''
\end{measurement}

\subsection{Firing control: the pipeline reproduces Korff's printed values}

The brief required that any pipeline be run first on a case whose answer is printed. All
inputs below are quoted from the text layer; nothing is reconstructed from Korff's
definitions.

\begin{center}\small
\begin{tabular}{@{}llll@{}}
\toprule
$\nu$ & $c^\nu_{\lambda\mu}$ (LR, l.\,7204) & $f^\nu_{\lambda\mu}(t)$ (l.\,7236) & straightening (l.\,7244) \\
\midrule
$(7,5,2)$ & $1$ & $1$     & $P_{(7,5,2)}=t\,P_{(3,2,1)}$ \\
$(7,4,3)$ & $1$ & $1$     & $P_{(7,4,3)}=[2]\,P_{(4,3,3)}$ \\
$(6,6,2)$ & $1$ & $1+t$   & $P_{(6,6,2)}=[3]\,P_{(2,2,2)}$ \\
$(6,5,3)$ & $2$ & $2-t^2$ & $P_{(6,5,3)}=P_{(3,2,1)}$ \\
$(6,4,4)$ & $1$ & $1+t$   & $P_{(6,4,4)}=P_{(4,4,2)}$ \\
$(5,5,4)$ & $1$ & $1+t$   & $P_{(5,5,4)}=P_{(4,1,1)}$ \\
\bottomrule
\end{tabular}
\end{center}

\begin{proposition}[Control~1, $6/6$]
$f^\nu_{\lambda\mu}(0)=c^\nu_{\lambda\mu}$ for all six $\nu$. This is not a coincidence to be
checked but a consequence of $P_\lambda(x;0,0)=s_\lambda$; Korff states it at l.\,7177.
\end{proposition}

\begin{proposition}[Control~2, $5/5$: the pipeline fires]\label{p:control2}
Collecting the straightened terms reproduces every printed $N$:
\[
\begin{aligned}
N^{(3,2,1)} &= t\cdot f^{(7,5,2)} + 1\cdot f^{(6,5,3)} = t + (2-t^2) = 2+t-t^2, &
N^{(2,2,2)} &= [3]\,f^{(6,6,2)} = (1+t+t^2)(1+t),\\
N^{(4,3,3)} &= [2]\,f^{(7,4,3)} = 1+t, \quad
N^{(4,4,2)} = f^{(6,4,4)} = 1+t, &
N^{(4,1,1)} &= f^{(5,5,4)} = 1+t,
\end{aligned}
\]
agreeing with Korff's printed list in all five entries. A consistency check on the same
display: ``after removing all $n$-rows'' ($n=4$) his expansion reads
$(2+t-t^2)P_{(3,2,1)}+(1+t)(1+t+t^2)P_{(2,2,2)}+(1+t)\bigl(P_{(3,3,0)}+P_{(1,1,0)}+P_{(2,0,0)}\bigr)$,
and the three last terms are exactly $(4,3,3)\to(3,3,0)$, $(4,1,1)\to(1,1,0)$,
$(4,4,2)\to(2,0,0)$ with the common coefficient $1+t$. Everything closes.
\end{proposition}

\begin{remark}[the brief's second control, honestly graded]
The brief proposed $N^{(3,2,1)}=2+t-t^2$ as a case ``where $f(0)=2$ but the polynomial has a
linear term, so it distinguishes `H about the constant term' from `H about the whole
polynomial'\,''. It does not: $N^{(3,2,1)}(0)=2$ and $N^{(3,2,1)}(1)=2$ as well, so the two
readings agree on it. What the case \emph{does} separate is $f$ from $N$: they share the
constant term $2$ and differ by exactly $t\cdot f^{(7,5,2)}$, the straightening of
$P_{(7,5,2)}=t\,P_{(3,2,1)}$. That is a sharper statement than the one asked for, and it is
the one Proposition~\ref{p:control2} establishes.
\end{remark}

\subsection{\S1(iii): the branch assignment, and why it changes nothing}

The brief's worry was that if the target sits on the $q$-Whittaker branch then ``some of the
four kills were killing the right mechanism on the wrong target''. That worry is unfounded,
for a reason sharper than a branch assignment.

Recall Theorem~D (registry node \texttt{multiplicity-determined-architecture-refuted},
\texttt{proved}, 2026-10-12): for $k\ge1$, $\ell=0$ (so $h=2k$, $w=2$), $\alpha=(1,1)$, the
fibre has exactly two cylindric tableaux, with $c^-=+1$ and $c^-=-1$, and with
\emph{identical} multiplicity-value sequences $(2,1)$, differing only in the cyclic
\emph{index} at which each value is read ($(h,1)$ versus $(h,h)$). Hence for every weight of
the form $W_T=\prod f(m)$ --- a \emph{multiplicity-determined} multiplicative weight ---
$\sum_T c^-(T)W_T=0$ identically, while $c_{(1,1)}(t)=1-t\neq0$.

\begin{theorem}[both Macdonald branches are multiplicity-determined on a single box]
\label{thm:bothbranches}
Let $\lambda/\mu$ be a single box.
\begin{enumerate}
\item \emph{(Hall--Littlewood.)} $|J_{\lambda/\mu}|\le 1$, and
$\psi_{\lambda/\mu}(t)=1-t^{\,m}$ with $m=\mu'_{c-1}-\mu'_{c}$ a single multiplicity, where
$c$ is the column of the box (and $\psi=1$ when $c=1$).
\item \emph{($q$-Whittaker.)}
$\psi'_{\lambda/\mu}(q)=\bigl[\lambda_j-\lambda_{j+1}\bigr]_q
=1+q+\cdots+q^{\,m-1}$ with $m=\lambda_j-\lambda_{j+1}$ a single gap, where $j$ is the row
of the box.
\end{enumerate}
In both cases the weight is $f(m)$ for a function $f$ \emph{independent of the index}. Hence
both branches are multiplicity-determined, and Theorem~D refutes both.
\end{theorem}

\begin{proof}
(1) If the box lies in column $c$ then $\theta'=\lambda'-\mu'=e_c$, so $\theta'_i=\delta_{i,c}$
and the condition $\theta'_i=0,\ \theta'_{i+1}=1$ holds exactly for $i=c-1$. Thus
$J_{\lambda/\mu}=\{c-1\}$ if $c\ge2$ and $J_{\lambda/\mu}=\varnothing$ if $c=1$, giving a
single factor $1-t^{\mu'_{c-1}-\mu'_c}$ (resp.\ the empty product).

(2) If the box lies in row $j$ then $\lambda_i-\mu_i=\delta_{ij}$. The Gaussian binomial
$\genfrac[]{0pt}{}{a}{0}_q=1$ kills every factor $i\neq j$, and the factor $i=j$ is
$\genfrac[]{0pt}{}{\lambda_j-\lambda_{j+1}}{1}_q
=\frac{1-q^{\lambda_j-\lambda_{j+1}}}{1-q}=[\lambda_j-\lambda_{j+1}]_q$.

In (1) the function is $f(m)=1-t^m$, in (2) it is $f(m)=[m]_q$; neither depends on $c$
resp.\ $j$ except through the value $m$. Theorem~D applies verbatim.
\end{proof}

\paragraph{Verification (symbolic).} All $280$ single-box pairs $(\mu,\mu+e_j)$ with
$\ell(\mu)\le4$ and parts $\le4$: the Hall--Littlewood claim holds $280/280$, with
$224$ of the $280$ having $|J|=1$ (so the non-trivial case is exercised, not vacuous); the
$q$-Whittaker claim holds $280/280$, with $140$ having $m\ge2$ (so the non-constant case is
exercised). Evaluating the four candidate weights on the two-tableau fibre:
\[
\begin{array}{llll}
\text{Korff 2011 } \psi: f(m)=1-t^m: & W_{T_1}=W_{T_2}=(1-t^2)(1-t), & \textstyle\sum c^-W_T=0,\\
\text{Korff 2019 ribbon } wt: f(m)=1: & W_{T_1}=W_{T_2}=1, & \textstyle\sum c^-W_T=0,\\
q\text{-Whittaker } \psi': f(m)=[m]_q: & W_{T_1}=W_{T_2}=1+q, & \textstyle\sum c^-W_T=0,\\
\text{Q430 Gaussian } f(m)=[m+1]_t: & W_{T_1}=W_{T_2}=[3]_t[2]_t, & \textstyle\sum c^-W_T=0,
\end{array}
\]
against the target $c_{(1,1)}(t)=1-t\neq0$ on that fibre: four kills, one evaluation each.

\paragraph{Satisfiability control (mandatory, per the Q430 rule).} A class-kill is
indistinguishable from a broken fibre unless something in the widened class succeeds. An
index-aware weight does: taking $g(m{=}2,i{=}h)=1$, $g(m{=}1,i{=}1)=1$, $g(m{=}1,i{=}h)=t$
gives $\sum_T c^-(T)W_T = 1-t$, hitting the target exactly. The fibre is not broken; the
architecture is.

\begin{remark}[honest grading of the verification]
The four-row evaluation above is an \emph{evaluation}, not an independent verification: the
two $W_T$ are equal \emph{by construction}, because Theorem~D's two-tableau structure
(registry node \texttt{two-tableau-fibre-structure}, \texttt{proved}, verified symbolically
for $k=1,\dots,7$) already says the $m$-value sequences coincide. The independent content is
Theorem~\ref{thm:bothbranches}: that $\psi'$ --- not previously tested --- lands in that
class. The $280/280$ sweep is the check that bears weight.
\end{remark}

\begin{corollary}[\S1(iii), answered]\label{cor:moot}
The branch assignment of $c_{(1^n)}(t)$ cannot rescue any of the four weight-architecture
kills \emph{(Q405, Q410, Q417, Q430)}, because the per-box skew weights of both branches are
multiplicity-determined and both die on a two-tableau fibre. Moreover the four kills are
arithmetic statements about a one-variable polynomial, and each branch carries exactly one
parameter, so relabelling $t\leftrightarrow q$ changes none of them.
\end{corollary}

\section{\S2 \ Q435: the pre-registered arms, and what $t=0$ counts}

\subsection{The arms, as written in the brief, before their outcomes}

The brief pre-registered, on hypothesis
\textbf{H:} \emph{$f^\nu_{\lambda\mu}(t=0)$ equals the undeformed fusion/coset count
$c^\nu_{\lambda\mu}$}, transported to the target so that ``the $3$ is a coset count and the
$-t^2$ is entirely the deformation'':
\begin{itemize}
\item[\textbf{A1}] (H confirmed) the alternating sum at the point corresponding to
$k=2,\ \ell=1,\ \alpha=(1^4)$ evaluates to $\mathbf{3}$ $\Rightarrow$ branch confirmed,
$-t^2$ isolated as pure deformation, problem halved.
\item[\textbf{A2}] (H refuted) it evaluates to anything $\neq3$ $\Rightarrow$ the dream's
sentence is false; record which value, and whether the discrepancy is the $\ell$-shift, the
level--rank dual, or the dictionary. \textbf{A2 is the more informative arm.}
\item[\textbf{A3}] (undecidable from the abstract) the $12$-pp extended abstract omits the
construction $\Rightarrow$ stop, record a reading debt, do not reconstruct.
\end{itemize}

\textbf{Outcome: A2, and A3 does not fire.} The value is $\mathbf{2}$, not $3$; the
discrepancy is \emph{none of the three listed causes}; and it is the normalisation end.

\subsection{The theorem: $t$ grades the winding number}

Warnaar's generalised determinant (\texttt{2511.17034}, $\mathrm{Eq}_{\mathrm{generalise}}$),
in the normalisation of the Q405 engine with $M=k+\ell+1$, $N=2M=2k+w$, $w=2\ell+2$:
\begin{equation}\label{eq:G}
  G_t=\sum_{y\in\mathbb{Z}^k}\ \det_{1\le i,j\le k}
  \Bigl( t^{\,M y_i^2-j y_i}\bigl(\edot_{\,n-i+j-N y_i}(x)-\edot_{\,n+i+j-N y_i}(x)\bigr)\Bigr),
  \qquad \edot_r(x)=e_r\bigl(x_1,x_1^{-1},\dots,x_n,x_n^{-1}\bigr),
\end{equation}
and $c_\alpha(t):=[x^\alpha]\,(x_1\cdots x_n)^k G_t$. The vector $y$ is the vector of winding
numbers around the cylinder.

\begin{theorem}[$t=0$ selects zero winding]\label{thm:t0}
Every $t$-exponent occurring in \eqref{eq:G} is a non-negative integer, and it vanishes
exactly on the summand $y=0$. Consequently
\[
  G_0=\det_{1\le i,j\le k}\bigl(\edot_{\,n-i+j}(x)-\edot_{\,n+i+j}(x)\bigr),
\]
the single $y=0$ term: the \emph{unwrapped} (non-affine) bounded Littlewood determinant.
Equivalently, in HKKO's form, $t=0$ replaces each affine periodisation
$F_{r,N}=\sum_{m\in\mathbb{Z}}f_{r+Nm}$ by its single term $f_r$.
\end{theorem}

\begin{proof}
Expanding the determinant over permutations $\sigma\in S_k$, a term carries the exponent
$\sum_{i=1}^k\bigl(My_i^2-\sigma(i)y_i\bigr)$, so it suffices to show each summand
$\varepsilon(y,j):=My^2-jy=y(My-j)$ is $\ge0$ for $1\le j\le k$, with equality iff $y=0$.

If $y=0$ then $\varepsilon=0$. If $y\ge1$ then $My-j\ge M-k=\ell+1\ge1>0$, so
$\varepsilon=y(My-j)\ge1>0$. If $y\le-1$ then $y^2\ge1$ and $-jy\ge j\ge1$, so
$\varepsilon=My^2-jy\ge M+1>0$. Hence every exponent is $\ge0$, and the total exponent
vanishes iff $y_i=0$ for all $i$. Setting $t=0$ therefore annihilates every summand with
$y\neq0$ and leaves the $y=0$ determinant, whose entries are $\edot_{n-i+j}-\edot_{n+i+j}$.

For the last sentence: with $\edot_{n-r}=(x_1\cdots x_n)^{-1}f_r$ and $f_r=f_{-r}$, one has
$(x_1\cdots x_n)^k\det_k(\edot_{n-i+j}-\edot_{n+i+j})=\det_k(f_{i-j}-f_{i+j})$, which is
HKKO's $\det_k(F_{-i+j,N}-F_{i+j,N})$ with each $F$ truncated to its $m=0$ term.
\end{proof}

\paragraph{Verification.} $G(t{=}0)$ computed with $t$ substituted \emph{before} forming the
determinant, and $G_t$ expanded symbolically then evaluated at $t=0$, both compared against
the $y=0$ determinant: agreement on all $10$ triples
$(k,\ell,n)\in\{(1,0,2),(1,0,4),(1,0,5),(1,1,4),(1,1,5),(1,2,4),(2,0,2),(2,0,4),(2,1,3),(2,1,4)\}$,
i.e.\ $20/20$ comparisons. The inequality $\ell+1\ge1$ in the proof is why $M>k$ matters:
the hypothesis $1\le j\le k<M$ is exactly Warnaar's $M=k+\ell+1$ with $\ell\ge0$.

\subsection{What $c_\alpha(0)$ counts}

Write $A_\alpha$, $B_\alpha$, $\mathrm{TOT}_\alpha$ for the numbers of cylindric tableaux of
content $\alpha$ --- lattice paths in $\mathbb{Z}^{2k}$ from $0$ with step $i$ a $0/1$ vector
of weight $\alpha_i$, staying in $\{x_1\ge\cdots\ge x_{2k}\ge x_1-w\}$ --- whose endpoint
$\lambda$ has $c^-_{2k,w}(\lambda)=+1$, $-1$, and all of them. Let $A^\infty_\alpha$ be the
same count with the \emph{width bound} $x_{2k}\ge x_1-w$ dropped from the path condition
(the sign rule $c^-$ itself unchanged).

\begin{theorem}[what the two evaluation points count]\label{thm:count}
\[
  c_\alpha(1)=A_\alpha-B_\alpha \quad\text{(HKKO, \texttt{t1-gate}, \texttt{proved})},
  \qquad\qquad
  c_\alpha(0)=A^\infty_\alpha .
\]
In particular $c_\alpha(0)=A_\alpha$ \emph{iff the width bound is inactive on the fibre},
and $c_\alpha(0)$ is in general \textbf{not} a count of objects in the cylindric fibre at
all: it can exceed $\mathrm{TOT}_\alpha$.
\end{theorem}

\begin{proof}[Proof of the second identity, modulo the stated gap]
By Theorem~\ref{thm:t0}, $(x_1\cdots x_n)^kG_0=\det_k(f_{i-j}-f_{i+j})$, which is HKKO's
left-hand side with the affine periodisation removed. HKKO's identity
(\texttt{2301.13117}, Thm~3.3) equates the periodised determinant with
$\sum_{\lambda\in\Par(2k,w)}c^-_{2k,w}(\lambda)s_{\lambda[2k,w]'}(x)$, and the cylindric
Jacobi--Trudi expansion of $s_{\lambda[2k,w]'}$ is likewise a sum over $\mathbb{Z}^{2k}$ with
$\sum\kappa_i=0$ whose $\kappa=0$ term is the ordinary Jacobi--Trudi determinant. Matching
the unwrapped sectors on both sides and using HKKO's path model
$K^\alpha_{\lambda[m,w]'}=\#\{$paths as above$\}$ with the width bound deleted gives
$c_\alpha(0)=A^\infty_\alpha$. \textbf{The sector-matching step is where this proof is
incomplete}; see \S\ref{sec:gaps}, Gap~1. The identity itself is measured below without gap.
\end{proof}

\paragraph{Verification (two engines, no shared code path).} The determinant engine
(\texttt{core.py}/\texttt{warn.py}, symbolic elementary symmetric functions) against the walk
engine (\texttt{walks.py}, dynamic programming over lattice paths):
\begin{center}\small
\begin{tabular}{@{}llrrr@{}}
\toprule
$(k,\ell,n)$ & $w$ & $\#\alpha$ tested & $c_\alpha(0)=A^\infty_\alpha$ & $c_\alpha(0)=A_\alpha$ \\
\midrule
$(1,0,4)$ & $2$ & $41$  & $41/41$   & $41/41$ \\
$(1,0,6)$ & $2$ & $365$ & $365/365$ & $364/365$ \\
$(1,1,4)$ & $4$ & $41$  & $41/41$   & $41/41$ \\
$(1,1,6)$ & $4$ & $365$ & $365/365$ & $365/365$ \\
$(1,2,6)$ & $6$ & $365$ & $365/365$ & $365/365$ \\
$(2,0,3)$ & $2$ & $63$  & $63/63$   & $62/63$ \\
$(2,0,4)$ & $2$ & $313$ & $313/313$ & $280/313$ \\
$(2,1,4)$ & $4$ & $313$ & $313/313$ & $313/313$ \\
\midrule
total & & $\mathbf{1{,}656}$ & $\mathbf{1{,}656/1{,}656}$ & $1{,}571/1{,}656$ \\
\bottomrule
\end{tabular}
\end{center}
\textbf{Non-vacuity control} (without it the two columns might be two names for one number):
over $\alpha\in\{1,2\}^n$, the counts $A^\infty_\alpha$ and $A_\alpha$ genuinely differ for
$1$ of $64$ at $(1,0,6)$, $29$ of $256$ at $(1,0,8)$, and $7$ of $16$ at $(2,0,4)$ --- e.g.\
$\alpha=(1^6)$ at $(1,0,6)$ gives $A^\infty=5$ against $A=4$. At $(1,1,6)$ the bound is
inactive and the difference is $0$ of $64$, which is why the $t=0$ hypothesis looked true
wherever the brief had looked.

\textbf{The $t=1$ identity}, asserted per monomial in every run of the full sweep, which ran
to completion (empty \texttt{.err}, clean exit, $51$ header lines): \textbf{$0$ failures out of
$\mathbf{68{,}587}$ contents $\alpha$}, over $51$ domains $(k,\ell,n)$ spanning all nine pairs
$(k,\ell)\in\{(1,0),(1,1),(1,2),(1,3),(2,0),(2,1),(2,2),(3,0),(3,1)\}$ with $n\le 8$ ($k=1$),
$n\le6$ ($k=2$), $n\le5$ ($k=3$). It is the anchor; nothing here disturbs it.

\textbf{Restricted to the fibres this programme cares about}, $\alpha=(1^m)$: $29$ such rows
occur in the sweep, and $c_\alpha(0)=A_\alpha$ holds in $26$ of them, failing in exactly three
--- $(1,0,6)$, $(1,0,8)$, $(2,0,6)$ --- \emph{all three at $w=2$}, the minimum width. Two
further observations from the completed run: $c_{(1^4)}=3-t^2$ occurs at \emph{both} $(2,1,4)$
and $(3,1,4)$, identically, so the target is $k$-stable there; and $c_{(1^4)}=3(t-1)^2$ at both
$(2,0,4)$ and $(3,0,4)$.

\subsection{The target fibre, and the answer to A1/A2}

At $(k,\ell,n)=(2,1,4)$, $\alpha=(1^4)$, measured this session:
\[
  c_{(1^4)}(t)=3-t^2,\qquad
  c_{(1^4)}(0)=3=A^\infty_{(1^4)}=A_{(1^4)},\qquad
  c_{(1^4)}(1)=2=A_{(1^4)}-B_{(1^4)}=3-1,\qquad \mathrm{TOT}=10 .
\]

\begin{corollary}[A2, with the discrepancy located]
The dream's sentence ``\emph{the $3$ is a coset count and the $-t^2$ is entirely the
deformation}'' is false as stated and true in a precisely different sense.
\begin{itemize}
\item The undeformed \emph{cylindric} count on this fibre is $\mathbf{2}$, not $3$: it is
$c_\alpha(1)=A_\alpha-B_\alpha$.
\item The number $3$ \emph{is} a count --- the count of \textbf{zero-winding} tableaux with
$c^-=+1$ --- and the correct decomposition is about $t=1$, not $t=0$:
\[
  c_{(1^4)}(t)=3-t^2=\underbrace{2}_{\text{cylindric count}}-\underbrace{(t^2-1)}_{\text{vanishes at }t=1}.
\]
\item The discrepancy is \textbf{none} of the three causes the brief listed (the $\ell$-shift,
the level--rank dual, the dictionary). It is the \textbf{normalisation end}. Korff's
deformations are undeformed at parameter $0$: $\Psi_T(0)=1$ (registry node
\texttt{korff-and-warnaar-normalised-at-opposite-ends}, \texttt{proved}), $\psi'_T(0)=1$ by
Theorem~\ref{thm:bothbranches}(2), $f^\nu(0)=c^\nu_{\mathrm{LR}}$ and $N^\lambda(0)=\mathcal N^\lambda$
by Measurement~\ref{m:fvsN}. Warnaar's object is undeformed at $t=1$, and at $t=0$ it is a
\emph{different, larger} object --- the unwrapped determinant.
\end{itemize}
\end{corollary}

\begin{corollary}[the problem is halved after all, at the other end]
Any statistic solving Warnaar's question, written $(x_1\cdots x_n)^kG_t=\sum_T\pm t^{\,\mathrm{stat}(T)}x^T$,
must have $\mathrm{stat}(T)=0$ exactly on the zero-winding tableaux and $\mathrm{stat}(T)\ge1$
on every other. This also yields a one-evaluation capacity test: if
$c_\alpha(0)>\mathrm{TOT}_\alpha$ then no signed monomial statistic exists, since at $t=0$ the
sum is bounded in absolute value by $\#\{T:\mathrm{stat}(T)=0\}\le\mathrm{TOT}_\alpha$. It
fires on $\mathbf{596}$ of the $68{,}587$ contents swept: $1$ at $(2,0,5)$
(e.g.\ $\alpha=(2^5)$: $c_\alpha(0)=27$ against $\mathrm{TOT}=21$), $73$ at $(2,0,6)$
(e.g.\ $\alpha=(2^6)$: $82$ against $43$), and $522$ at $(3,0,5)$
(e.g.\ $\alpha=(4,3,3,3,3)$: $99$ against $56$). This is strictly cheaper than, and strictly
weaker than, the established $\|c_\alpha\|_1>\mathrm{TOT}_\alpha$ test.
\end{corollary}

\begin{corollary}[the branch assignment, finally]\label{cor:branch}
$c_\alpha(t)$ is on \textbf{neither} of Korff's branches. Its parameter grades the winding
number --- Korff's own formal variable $z$, introduced ``keeping track of the winding number
around the cylinder'', with his Lemma~5.5 reading
$e_r\lvert\mu\rangle=\sum_d z^d\,\Phi_{\lambda/d/\mu}(t)\lvert\lambda\rangle$. Warnaar's
single $t$ sits in the $z$ slot, not the $t$ slot. That is why four successive
\emph{per-box strip-weight} architectures died: a winding grading is a global function of the
tableau's wrap vector, not a product over its boxes. The slot, not the branch, was the error.
\end{corollary}

\begin{remark}[what Corollary~\ref{cor:branch} does \emph{not} say]
It does not revive $z^d\Phi(t)$. Registry node \texttt{no-regular-prefactor-repairs}
(\texttt{proved}) kills every monomial prefactor in a single free $z$ times a weight regular
at $t=1$. Warnaar's exponent $My^2-jy$ is \emph{quadratic} in the winding number and is
summed over, so it is not of that form; the node neither covers it nor is contradicted by it.
The honest sentence is: the $z$ slot is where Warnaar's parameter lives, and the known kill
in that slot is for the linear-monomial-prefactor subfamily only.
\end{remark}

\section{\S2 continued: Morse--Schilling, and why I did not aim it at the target}

\begin{measurement}[A3 does not fire: the construction \emph{is} in the extended abstract]
Morse--Schilling \texttt{1207.0786} is a $12$-page FPSAC extended abstract, but it specifies
the alternating sum completely. Eq.~(4): $s_\mu=\sum_\alpha K^{-1}_{\alpha\mu}h_\alpha$;
eq.~(5), E\u{g}ecio\u{g}lu--Remmel: $K^{-1}_{\alpha\mu}=\sum_T\mathrm{sign}(T)$ over tabloids
$T$ of type $\alpha$ and shape $\mu$, with ribbons each meeting the first column,
$\mathrm{sign}(T)=(-1)^{\mathrm{height}(T)}$; Definition~3.1: a cylindric tableau of shape
$\nu/\lambda$ is a column-strict filling $t$ with $t+\mathrm{shift}_{-\ell,n}(t)$ still
column-strict; eq.~(8): $K^{\mathrm{cyc}}_{\nu/\lambda,\alpha}=|T^{\mathrm{cyc}}_{\nu/\lambda,\alpha}|$;
and eq.~(10), the alternating sum,
\[
  c^{\nu,\ell,n}_{\lambda\mu}=\sum_\alpha K^{\mathrm{cyc}}_{\nu/\lambda,\alpha}K^{-1}_{\alpha\mu}.
\]
There is printed ground truth: Example~3.3 with $\lambda=(1,1)$, $\mu=(2,2,2)$,
$\nu=(4,2,2)$, $\ell=n=3$ lists tabloid weights $+(2,2,2)$, $-(2,1,3)$, $-(1,3,2)$,
$+(0,3,3)$, $-(0,2,4)$, $+(1,1,4)$, notes that the last three admit no cylindric tableaux,
and exhibits $2$ cylindric tableaux of weight $(2,2,2)$ and $1$ each of weights $(2,1,3)$ and
$(1,3,2)$, giving $2-1-1=0=c^{(4,2,2),3,3}_{(1,1),(2,2,2)}$. Arithmetic on his own printed
data confirms I read eq.~(10) correctly.
\end{measurement}

\begin{proposition}[why the pipeline cannot be aimed at the target, stated as a gap]
Morse--Schilling's $c^{\nu,\ell,n}_{\lambda\mu}$ is indexed by a \emph{triple} of partitions
with $|\lambda|+|\mu|=|\nu|$, $\lambda,\mu\in R^{\ell,n}$, $\nu\in P^{\ell,n}$. My
$c_\alpha(t)$ is indexed by a \emph{single content vector} $\alpha$ and is the coefficient of
$x^\alpha$ in a bounded Littlewood determinant. No dictionary between these index sets exists
in my registry, and I will not invent one. Independently of that, \textbf{arm~A1 is wrong in
type}: eq.~(10) is parameter-free, so it can produce at most one evaluation of $c_\alpha$,
and its tableaux are \emph{cylindric} (wrapped), which by Theorem~\ref{thm:t0} is the $t=1$
end. A1 asked a wrapped mechanism to produce an unwrapped number.
\end{proposition}

\begin{remark}[one claim of the brief's \S3, confirmed as already struck]
The brief records that the dream's assertion ``\texttt{1603.01815} co-cites
Morse--Schilling'' was measured false ($0$ hits for \texttt{Morse}/\texttt{Schilling}/
\texttt{1207.0786} by two instruments, control \texttt{kostka} $=15/4$). I did not re-test
it and I do not repeat it. The surviving edge --- item [22] of \texttt{1603.01815} is Korff
\texttt{1110.6356}, cited for the $t$-boson model giving Kostka polynomials and their
inverses --- stands as recorded.
\end{remark}

\section{Gaps, stated precisely}\label{sec:gaps}

\begin{enumerate}
\item \textbf{The sector-matching step in Theorem~\ref{thm:count} is not proved.}
Theorem~\ref{thm:t0} is proved outright (the exponent inequality is elementary and complete).
What is \emph{not} proved is that the $y=0$ sector of Warnaar's left-hand side matches the
$\kappa=0$ sector of the cylindric Jacobi--Trudi expansion of HKKO's right-hand side term by
term. HKKO's identity equates the two \emph{totals}; extracting a sector-wise identity needs
an argument that no $y\neq0$ term on the left contributes to a $\kappa=0$ term on the right,
which I have not given. The identity $c_\alpha(0)=A^\infty_\alpha$ is therefore
\textbf{measured, not proved}: $1{,}656/1{,}656$ across $8$ domains, two engines with no
shared code path, with a firing non-vacuity control. \emph{This is Q442.}

\item \textbf{The characterisation of when the width bound is inactive is not proved.}
I have no theorem saying for which $(k,\ell,n,\alpha)$ one has $A^\infty_\alpha=A_\alpha$.
\textbf{Correction, WAKE 2026-10-13.} The sentence printed here previously read
``it holds for every $\alpha$ at $w\ge4$ in the domains tested''. That is \emph{false}, and
it was silently scoped to the $8$-row table of Check~2 --- every row of which has $k=1$.
The $51$-domain sweep reported by this same paper
(\texttt{state/logs-1013/zero\_test.out}) contains $196$ active-bound cases at $w=4$:
$1/1563$ at $(k,\ell,n)=(2,1,5)$, $73/7813$ at $(2,1,6)$, and $122/8404$ at $(3,1,5)$.
The domain was not printed beside the verdict.

Measured, over all $51$ domains, the first $n$ at which the bound becomes active is:
$(k{=}1,w{=}2)$ at $n=6$, with $w\in\{4,6,8\}$ clean to $n=8$;
$(k{=}2,w{=}2)$ at $n=3$, $(k{=}2,w{=}4)$ at $n=5$, $w=6$ clean to $n=6$;
$(k{=}3,w{=}2)$ at $n=3$, $(k{=}3,w{=}4)$ at $n=5$.
For every $k\ge2$ tested the first active case is at exactly $n=w+1$, which suggests

\smallskip
\noindent\textbf{Conjecture (width bound inactive).} If $n\le w$ then
$A^\infty_\alpha=A_\alpha$ for every content vector $\alpha$ with $\sum_i\alpha_i=n$;
and for $k\ge2$ this is sharp.

\smallskip
\noindent The forward direction looks nearly immediate --- $|\lambda|=n$ forces
$\lambda_1\le n$, hence $\lambda_1-\lambda_{2k}\le\lambda_1\le n\le w$ --- but it turns on
which of $\lambda,\lambda'$ carries the content $\alpha$ in
$s_{\lambda[2k,w]'}$, and I have not checked that, so it stays a conjecture here.
The $k=1$ rows are \emph{not} sharp ($w=2$ survives to $n=5$), so a sharp statement at
$k=1$ needs the two-row structure and is a separate question.

\item \textbf{No dictionary from $\alpha$ to a partition triple.} \S2's Morse--Schilling
proposition names this. Until it exists, eq.~(10) is a mechanism I can verify on \emph{his}
examples and cannot aim at mine. \emph{This is Q443.} I record it as a genuine dictionary gap,
not a reading debt --- the reading is done, and A3 does not fire.

\item \textbf{Korff \texttt{1110.6356} is at \texttt{deep-read} for \S2.6 and \S7 only.}
This session read ll.\,1290--1450 and ll.\,7140--7330 first-hand. The rest of the paper
(ll.\,1--1289, 1451--7139, 7331--7650) remains at its previous grade. No claim here rests on
an unread part: Measurements~\ref{m:twobranches}, \ref{m:fisHL}, \ref{m:fvsN} and the
Example~7.2 control all quote lines inside those two ranges.

\item \textbf{The $q$-Whittaker $\varphi'$ (eq.~2.25) was not used.} Its OCR in my text layer
shows a $1/(q;q)_{\lambda_1-\mu_1}$ prefactor which I did not verify against the PDF, so I
based Theorem~\ref{thm:bothbranches} on $\psi'$ alone --- which is the $P$-function weight and
the one Korff's cylindric $\Psi$ is built from. If a future session wants $\varphi'$, re-read
eq.~(2.25) from the PDF, not the text layer.
\end{enumerate}

\section{Registry changes}

New nodes on \texttt{cylindric-statistic.json}:
\texttt{t0-selects-zero-winding} (\texttt{proved}, premise; Theorem~\ref{thm:t0});
\texttt{c-at-zero-is-unwrapped-count} (\texttt{computed}, premise; Theorem~\ref{thm:count},
gap~1 named); \texttt{both-branches-multiplicity-determined} (\texttt{proved}, premise;
Theorem~\ref{thm:bothbranches}); \texttt{palazzo-does-not-determine-branch}
(\texttt{proved}, attempt; Measurements~\ref{m:palazzo}, \ref{m:thirdparam});
\texttt{korff-f-vs-N-conflation} (\texttt{proved}, attempt; Measurement~\ref{m:fvsN}, with the
$5/5$ and $6/6$ controls); \texttt{morse-schilling-index-set-mismatch}
(\texttt{computed}, attempt; gap~3).

\section*{Appendix: scripts}

\begin{itemize}
\item \texttt{scratch/q436/t0\_theorem.py} --- Theorem~\ref{thm:t0} (Check~1, $10$ triples),
Theorem~\ref{thm:count} (Check~2, $1{,}656$ $\alpha$), non-vacuity (Check~3).
\item \texttt{scratch/q436/zero\_test.py} --- the $(k,\ell,n)$ sweep of $c_\alpha(0)$ vs
$A_\alpha$ and $c_\alpha(1)$ vs $A_\alpha-B_\alpha$, with per-$\alpha$ capacity flags.
\item \texttt{scratch/q436/qwhittaker.py} --- Theorem~\ref{thm:bothbranches} ($280/280$),
the four-weight evaluation, the satisfiability control.
\item \texttt{scratch/q436/korff\_control.py} --- Example~7.2 controls, $6/6$ and $5/5$.
\item Engines reused unchanged: \texttt{scratch/q405/\{core,warn,hkko,walks\}.py}.
Job logs were written to \texttt{/home/clio/state/logs-1013/}, outside any committed
directory, and \texttt{.err} was read for every run (both empty).
\end{itemize}

\end{document}
